


\section{Experimental Evaluation}
\label{sec:experiments}

We implemented a prototype of \system\ and evaluated it on several datasets prepared for the \emph{Text-to-Constraints} setting. Our evaluation is organized around the following questions.
\begin{itemize}[leftmargin=0.5cm]
    \item \textbf{RQ1: Constraint-sentence classification.} How accurately can \system\ identify which sentences in a normative document express data constraints?

    \item \textbf{RQ2: Mapping quality.} Given constraint-relevant text, how accurately can \system\ map it to executable denial constraints and  how well do these constraints capture the corresponding data violations?

    \item \textbf{RQ3: Impact of implementation choices.} How do different classifier and mapper implementations, LLM backends, and prompt variants affect extraction quality and robustness?

    \item \textbf{RQ4: Comparison with baselines.} How does \system\ compare with data profiling and Text-to-SQL baselines on constraint extraction and violation detection?
\end{itemize}

\subsection{Datasets}
\label{sec:datasets}

%No benchmark currently exists for the \emph{Text-to-Constraints} problem.
%We prepared four datasets for our experiments. Two are based on real-world regulatory documents paired with corresponding tabular datasets, and two are derived from datasets widely used in data quality and profiling research, for which the target constraints are already known. For the real-world datasets, we manually annotated the constraints expressed in the source text. For the derived datasets, we generated natural-language descriptions of the known constraints, as described below. When a constraint had no natural violations in the data, we injected a few violating records so that every constraint could be evaluated under the violation-based metrics of Section~\ref{sec:metrics}. 
%Table~\ref{tab:datasets} summarizes the datasets.


We evaluate \system\ on six datasets: four based on real-world regulatory documents paired with corresponding tabular datasets, and two derived from benchmark datasets commonly used in data quality and profiling research. A summary of these datasets is provided in Table~\ref{tab:datasets}, including the number of rows and columns, the number of sentences in the source text, the number of annotated constraints, and the number of tuples violating at least one constraint.



\begin{table}[ht]
  \centering
\small
%  \resizebox{\columnwidth}{!}{%
  \begin{tabular}{lrrrrr}
    \toprule
    \textbf{Dataset} & \textbf{\#Rows} & \textbf{\#Cols} & \textbf{Sent.} & \textbf{Const.} & \textbf{Viol.} \\
    \midrule
    Hospital & 1k & 17 & 109 & 15 & 834 \\
    Restaurant & 128k & 21 & 33 & 7 & 52,311 \\
    SISU & 102k & 56 & 42 & 17 & 533 \\
    Tax & 100k & 15 & 126 & 6 & 468 \\
    HMDAZ$^{\dagger}$ & 1M & 135 & 17 & 9 & 209 \\
    HMDA Checks$^{\dagger}$ & 1M & 81 & 158 & 92 & 986,491 \\
    \bottomrule
  \end{tabular}%
  %}% resizebox
    \caption{Dataset statistics. $^{\dagger}$First 1M rows of the public file (15.1M and 4.8M rows); \emph{Const.} counts sentences mapped to a non-null gold DC.}
  \label{tab:datasets}
\end{table}

\smallskip\noindent\textbf{SISU.}
As discussed in Section~\ref{sec:intro}, the SISU dataset contains admission
records from Brazil's centralized university selection system. We manually
annotated three sections of \textit{Portaria Normativa No.~21/2012} covering student
registration, candidate classification and selection, and waitlist procedures.
The resulting constraints span nullability checks, functional dependencies,
ordering constraints, and conditional selection rules. Many have violations in
the released data, suggesting inconsistencies between the published records and
the regulatory requirements expressed in the text.
\altic{[R1-O10] Detalhar a construção do ground truth: (i) como as gold constraints foram derivadas/anotadas e por quem, (ii) similaridade e redundância entre sentenças (há sentenças quase duplicadas?), e (iii) classificar os tipos de constraint do SISU em formalismos conhecidos: ordering = order dependencies (ODs), conditional selection = conditional FDs (CFDs)? O revisor pediu esse mapeamento explícito.}

\smallskip\noindent\textbf{Restaurant.}
As introduced in Example~\ref{ex:classification}, this dataset contains
restaurant inspection records published by the Department of Health and Mental
Hygiene of New York. We annotated the NYC food service operator
inspection guide, focusing on the grading and closure sections, which present
constraints on score-to-grade mappings, inspection frequency, grade posting
timing, and closure conditions. Most have natural violations and can be
grounded directly to existing fields. For the one constraints that reasons over
inspection sequences, we added one derived column that tracks the days until the next inspection. We leave to future work how \system\ can guide or automate such dataset
transformations more systematically.


\smallskip\noindent\textbf{Tax and Hospital.}
Both datasets have been widely used in data profiling and denial-constraint
research since the seminal work of Chu et al.~\cite{chu2013discovering}. For
Tax, we extracted the constraint descriptions from the original
paper~\cite{chu2013discovering}. For Hospital, we used the data and constraints from
a well-known cleaning system~\cite{holoclean}. Because their constraints are available in
formal form, we first generated natural-language descriptions with an LLM. For
both datasets, we then prompted an LLM to generate additional non-constraint
sentences around each rule, mimicking a regulatory document that interleaves
binding rules with other prose. The prompt receives the schema and is
instructed to use schema terminology, making sentence classification more
challenging. To avoid circularity between data generation and evaluation, we used Claude~4.5 Sonnet, a model not
evaluated in any \system\ configuration.
\altic{[R3-O3, R2-O6.2] Assumir explicitamente a limitação destes dois datasets: o ``texto regulatório'' vem dos papers que já enunciam as regras, e as sentenças são geradas por LLM, compartilhando priors com os backbones avaliados. Isso infla os scores (Hospital F1=1.0) e superestima a dificuldade real. Acrescentar uma frase deixando claro que SISU e Restaurant (pares texto-tabela realistas, F1 perto de 0.5) são o sinal mais confiável, e que Tax/Hospital servem como cenário controlado.}



\smallskip\noindent\textbf{HMDAZ and HMDA Checks.}
Both datasets use public loan-level records released under the Home Mortgage
Disclosure Act (HMDA) and are, by far, the largest in our evaluation; we use the
first 1M rows of each file. \emph{HMDAZ} pairs the 2018 loan application
register (135 columns, including the code and label versions of each
categorical field) with Sections 1026.32, 1026.35, and 1026.43 of Regulation~Z,
which define high-cost (HOEPA) and higher-priced mortgages, the limitations
that apply to them (e.g., balloon payments, negative amortization, prepayment
penalties), and the ability-to-repay requirements. Sentences are verbatim from
the regulation, and gold DCs keep the coverage conditions of each trigger (e.g.,
a loan is high-cost only if it is a closed-end, principal-residence loan).
Scope-only text, or text whose outcome is not observable in the data, is
annotated with a \texttt{null} constraint. \emph{HMDA Checks} pairs the 2014
register (81 columns) with the 2014 HMDA edit specification published by the
FFIEC: one entry per edit, with the edit test verbatim; edits that cannot be
checked on a single record (aggregate, transmittal-sheet, or external-reference
edits) are \texttt{null}. These datasets stress three aspects that the original
four do not: wide schemas whose relevant column is often named after a code
rather than the concept in the text, rules that compare columns with each other
or through arithmetic (e.g., loan amount versus five times the income), and a
highly concentrated violation distribution (Section~\ref{sec:metrics}).

\subsection{Baselines}
\label{sec:baselines}

Since no prior system addresses the Text-to-Constraints task directly, we compare \system\ against the two closest paradigms: DC discovery from data alone and Text-to-SQL.
\altic{[R1-O12] Considerar baselines além de DC discovery, já que reivindicamos cobrir ordering e conditional selection: incluir descoberta de ODs (order dependencies) e CFDs como comparativos, ou justificar por que DC discovery aproximada já subsume essas classes no nosso setting. Conectar com a classificação de tipos de constraint pedida em [R1-O10].}

\smallskip\noindent\textbf{Data profiling.}
We use the DC discovery algorithms from~\cite{Pena2022FastDC}, running both exact and approximate denial-constraint discovery. Exact discovery returns only constraints that hold on the data with no violations. As a result, it is not suitable for our violation-based evaluation, which aims to recover rules that may be violated in the released dataset. We therefore focus on approximate DC discovery as the main profiling baseline. Approximate DCs allow a small fraction of violating tuple pairs, controlled by a threshold $\varepsilon$, and can therefore recover rules even when the data is not fully clean.

Because discovery can return many DCs, we rank them using the standard
criterion from prior work~\cite{chu2013discovering}, which combines coverage
and succinctness, with succinctness measured as the inverse of the number of
symbols in the DC. We run a grid over multiple top-$k$ values and
$\varepsilon$ thresholds, and report the best results obtained across these settings, following standard practice in prior work.

These baselines do not scale to our largest schemas. In particular, they
cannot process SISU or Restaurant with all columns present. For those two
datasets, we therefore restrict the input to the subset of columns that appear
in the gold constraints, giving the profiling baselines the best possible
setting under their scalability limits.


\smallskip\noindent\textbf{Text-to-SQL.}
We use XiYan-SQL~\cite{XiYanSQL}, a state-of-the-art generative text-to-SQL
framework. We adapt it to our setting by prompting it, for each sentence, to
return a SQL query that identifies the tuples violating the described
constraint, or \texttt{NULL} if the sentence does not express a checkable data
constraint. Under this formulation, XiYan-SQL acts jointly as a constraint
classifier and mapper, making it directly comparable to \system\ on  across all
evaluation dimensions.


\subsection{Evaluation Metrics}
\label{sec:metrics}

Text-to-Constraints requires solving three related problems: identify
text that expresses constraints, translate those into a formal
representation, and identify the data violations implied by those constraints.
We therefore evaluate systems along three complementary dimensions:
\textit{Constraint Classification}, \textit{Constraint Extraction}, and
\textit{Violation Detection}. This is similar in spirit to Text-to-SQL evaluation, where
execution-based evaluation complements structural matching because exact form
alone does not fully capture correctness~\cite{zhong-etal-2020-semantic}. Because \system\ and the
baselines in Section~\ref{sec:baselines} differ in what they produce, not all
dimensions apply to every method.

\smallskip\noindent\textbf{Constraint Classification.}
This dimension measures whether a system correctly identifies which sentences
in a document express constraints. We report precision, recall, and F$_1$
over all sentences. This dimension applies to sentence-level variants of
\system\ and to the Text-to-SQL baseline. It does not apply to data profiling
baselines, which operate on data rather than text.

\smallskip\noindent\textbf{Constraint Extraction.}
This dimension measures how accurately systems translate text into denial
constraints. We compare predicted and ground-truth constraints by
predicate-set match (i.e., exact matching over the set of predicates), and report precision, recall, and F$_1$. For \system\
and the data profiling baselines, the predicate set is read directly from the
predicted denial constraint. For Text-to-SQL, we extract predicates from the
\texttt{WHERE} clause of the generated query. Because different SQL
formulations may express the same underlying constraint, this metric can
underestimate Text-to-SQL performance, which further motivates reporting the
other evaluation dimensions. For data profiling, discovered constraints are
matched directly against the ground-truth set.

\smallskip\noindent\textbf{Violation Detection.}
This dimension measures whether the constraints produced by a system, when
executed on the dataset, identify the same tuples as the
ground-truth constraints. Because a denial constraint may involve either a
single tuple or a pair of tuples, we evaluate violations at the tuple level by
counting tuples that participate in at least one violation. We execute all predicted
constraints and all ground-truth constraints on the dataset, then compare the
resulting violation sets using precision, recall, and F$_1$. This dimension
applies to all \system\ variants and all baselines.


\smallskip\noindent\textbf{Complementary mapping metrics.}
Exact predicate-set matching penalizes logically equivalent rewritings, and the
tuple-level aggregate can be dominated by a single rule. We therefore report
four complementary measures for the mappers. (i)~\emph{Equivalence-aware
extraction} (CE$_{\equiv}$) compares DCs after a logical normal form that expands
\texttt{IN}/\texttt{NOT IN}, orients and sorts operands, flattens and
distributes \texttt{AND}/\texttt{OR}, and removes absorbed terms; in addition,
an equality literal on a code column and the equality literal on its label
column (e.g., \texttt{hoepa\_status = 1} and \texttt{hoepa\_status\_name = 'High-cost
mortgage'}) are unified only when the data confirms that both select exactly the
same tuples. Code--label pairs are discovered in the data, not from column
names. (ii)~\emph{Execution equivalence} is the fraction of gold rules with at
least one violation whose predicted DC flags exactly the same tuples.
(iii)~\emph{Per-rule violation F$_1$} averages, over gold sentences, the F$_1$
between the tuples flagged by the gold and the predicted DC, so that every rule
weighs the same. (iv)~\emph{Violation concentration} is the share of the
gold-violating tuples flagged by the single largest gold rule; it is a property
of the dataset, not of the method. We also count predicted DCs whose SQL fails
to execute, instead of treating them as producing no violation.


\subsection{Implementation and Execution Settings}
\label{sec:implementation}

\system\ is implemented in Python, with DuckDB as the query execution backend.
All LLM calls are issued through OpenRouter with strict JSON-schema structured
outputs. We evaluate six backbones: GPT-6 Luna, DeepSeek-V4-Flash, Mistral Large 4,
Qwen3.8-2.4T-A95B, GPT-6 Sol, and Claude Sonnet 5.5. Before the runs, every
backbone passed a health check of the classifier, both single-prompt output
schemas, and each Multi-Step schema. Two candidates were replaced: Kimi K2.6,
which repeatedly emitted whitespace until the token limit under strict
JSON-schema decoding, and MiniMax-M3, which returned empty single-prompt outputs
in the experiments (Section~\ref{sec:constraintextraction}). Each configuration is executed once, with a cap of 32k
output tokens per call. Unless stated otherwise, \system\ uses the schema-driven
predicate space of Section~\ref{sec:mapper:predicate-space} with on-demand predicate
generation, and the single-prompt mapper uses semantic column retrieval and
column context (Section~\ref{sec:mapper:llm}). Constraints are executed on the
full datasets (first 1M rows for HMDA) with DuckDB. \system\ experiments were
executed on a Linux workstation with an Intel Core i7-7700HQ CPU and 16\,GB of
memory; the DC discovery baselines were executed on a MacBook
Pro with an Apple M1 chip and 16\,GB of memory. The XiYan-SQL
baseline, which requires GPU inference, was executed on a server with an
NVIDIA A100 SXM GPU (80\,GB VRAM), 125\,GB RAM, and 16\,vCPUs.

%\smallskip\noindent\textbf{Efficiency} reports total wall-clock runtime and, for
%LLM-based approaches, estimated API cost per dataset, enabling analysis of the
%quality-cost trade-off across \textsc{Norma} configurations and baselines.



% \subsection{Implementation and Execution Settings}
% \label{sec:implementation}

% \system\ is implemented in Python, using DuckDB as the query execution
% backend. LLM calls are
% issued via the OpenAI API (GPT-5, GPT-5-mini), the Mistral API
% (Mistral-Small, Mistral-Large), the Groq API (Llama-3.3-70B), and the Qwen API (Qwen3-32B).
% All \system\ experiments and DC discovery baselines were executed on
% a MacBook Pro with an Apple M1 chip and 16\,GB of unified memory.
% Baselines requiring GPU inference were executed on a cloud server
% equipped with an NVIDIA A100 SXM GPU (80\,GB VRAM, 125\,GB RAM,
% 16\,vCPUs).
